\documentclass[11pt,a4paper]{report}

\usepackage[utf8]{inputenc}
\usepackage[T1]{fontenc}
\usepackage[spanish,es-tabla]{babel}
\usepackage{graphicx}
\usepackage{booktabs}
\usepackage{tabularx}
\usepackage{float}
\usepackage{geometry}
\usepackage{xcolor}
\usepackage{listings}
\usepackage[hidelinks]{hyperref}

\geometry{margin=2.5cm}

\lstset{
  language=bash,
  basicstyle=\ttfamily\footnotesize\color[RGB]{229,229,229},
  breaklines=true,
  breakatwhitespace=false,
  columns=fullflexible,
  keepspaces=true,
  showstringspaces=false,
  backgroundcolor=\color{black},
  frame=none,
  keywordstyle=\color[RGB]{86,182,194},
  commentstyle=\color[RGB]{106,153,85},
  stringstyle=\color[RGB]{206,145,120},
  identifierstyle=\color[RGB]{229,229,229},
  numberstyle=\color[RGB]{127,127,127},
  tabsize=2
}

% ------------------------------------------------------------------
% Macro de figura tolerante: si la imagen todavia no existe, dibuja
% un recuadro con el nombre del archivo que falta (asi el .tex compila
% mientras se van agregando las capturas a img/lab04/).
%   \fig[ancho]{archivo.png}{epigrafe}
% ------------------------------------------------------------------
\newcommand{\fig}[3][0.92]{%
  \begin{figure}[H]
    \centering
    \IfFileExists{img/lab04/#2}{\includegraphics[width=#1\linewidth]{img/lab04/#2}}%
      {\fbox{\parbox[c][4.5cm][c]{0.85\linewidth}{\centering\itshape [FALTA la captura: \texttt{img/lab04/#2}]}}}
    \caption{#3}
  \end{figure}
}

\begin{document}

\sloppy

\begin{titlepage}
  \centering
  \vspace*{1cm}

  {\large Laboratorio No. 04 --- Application Layer and Physical Layer Protocols \par}
  \vspace{1.5cm}

  {\large Por: \par}
  {\large Camilo Aguirre \par}
  {\large Juan David Rangel \par}
  \vspace{1cm}

  {\large Profesor: John Alexander Pachón Pinzón \par}
  \vspace{1.5cm}

  {\large Escuela Colombiana de Ingeniería Julio Garavito \par}
  {\large Universidad \par}
  \vspace{0.5cm}

  {\large ISIS AYSR - 1 \par}
  {\large Arquitectura y Servicios de Red \par}
  \vspace{0.5cm}

  {\large Asignatura: Arquitectura y Servicios de Red (AYSR) \par}
  \vspace{1cm}

  {\large Grupo: 2 estudiantes --- Camilo Aguirre, Juan David Rangel \par}
  {\large Fecha: \today \par}
  \vspace{1.5cm}

  {\large Bogotá, D.C., \today \par}

  \vfill
\end{titlepage}

\tableofcontents
\newpage

\section{Introducción}

En este laboratorio se prueban y analizan los \textbf{protocolos de capa de aplicación} de la
infraestructura, y se realizan actividades de \textbf{capa física} (cableado estructurado).
El trabajo tiene dos grandes escenarios:

\begin{enumerate}
  \item Una \textbf{red simulada en Cisco Packet Tracer}, con tres sedes
        (\texttt{sistemas.com}, \texttt{civil.com} y \texttt{electrica.com}) que ofrecen
        servicios de \textbf{DNS, HTTP, correo (SMTP/POP3) y FTP}.
  \item La \textbf{red real}, donde se monitorean protocolos con \textbf{Wireshark}, se
        consulta información de dominios en \textbf{centralops.net}, se monta un
        \textbf{servidor NTP} y se construyen \textbf{cables de red} (patch cords y
        crimpado en patch panel).
\end{enumerate}

\section{Marco teórico}

\subsection{Protocolos de capa de aplicación y sus puertos}

Los protocolos de la capa de aplicación se apoyan en TCP o UDP y se identifican por su
\textbf{puerto} en la capa de transporte.

\begin{table}[H]
  \centering
  \begin{tabularx}{\linewidth}{@{}llX@{}}
    \toprule
    \textbf{Protocolo} & \textbf{Puerto} & \textbf{Función} \\
    \midrule
    DNS  & 53 (UDP/TCP)  & Traduce nombres a direcciones IP. \\
    HTTP & 80 (TCP)      & Transferencia de páginas web. \\
    HTTPS & 443 (TCP)    & HTTP cifrado. \\
    SMTP & 25 (TCP)      & \emph{Envío} de correo. \\
    POP3 & 110 (TCP)     & \emph{Recepción/descarga} de correo. \\
    IMAP & 143 (TCP)     & Recepción/consulta de correo en el servidor. \\
    FTP  & 21 y 20 (TCP) & Transferencia de archivos (control y datos). \\
    DHCP & 67/68 (UDP)   & Asignación automática de direcciones IP. \\
    \bottomrule
  \end{tabularx}
  \caption{Puertos de los protocolos usados en el laboratorio.}
\end{table}

\subsection{Cableado estructurado}

El cableado estructurado organiza el cableado físico de un edificio en subsistemas
(entrada, backbone, horizontal, área de trabajo) usando conectores estándar
(\textbf{RJ-45}) y organización central (\textbf{patch panel}). El estándar define dos
asignaciones de color para los 8 hilos del par trenzado, \textbf{T568A} y \textbf{T568B};
un cable \emph{directo} usa la misma norma de los dos extremos, y uno \emph{cruzado}
usa T568A en un extremo y T568B en el otro.

\subsection{NTP}

\textbf{NTP} (Network Time Protocol, puerto 123/UDP) sincroniza el reloj de los equipos de
una red. Es importante porque muchos servicios (logs, autenticación, correo, bases de
datos, certificados) dependen de una hora consistente.

\section{Desarrollo}

\subsection{Parte A --- Cisco Packet Tracer}

\subsubsection{Topología}

La red tiene tres sedes, cada una con su switch (2950), en topología de estrella: un
switch central (\texttt{Switch0}) conecta a los otros dos. Se usa máscara
\textbf{255.0.0.0 (/8)}, por lo que todos los dispositivos quedan en la misma red
\texttt{27.0.0.0/8} y \textbf{no se requiere router}.

\fig{fig-pt-topologia.png}{Topología de la red de las tres sedes en Packet Tracer.}

\subsubsection{Direccionamiento}

\begin{table}[H]
  \centering
  \begin{tabularx}{\linewidth}{@{}lllX@{}}
    \toprule
    \textbf{Sede} & \textbf{Dispositivo} & \textbf{IP} & \textbf{Gateway / DNS} \\
    \midrule
    sistemas  & www.sistemas.com  & 27.100.0.32  & 27.0.0.1 / 27.1.0.2 \\
    sistemas  & mail.sistemas.com & 27.100.0.33  & 27.0.0.1 / 27.1.0.2 \\
    sistemas  & DNS               & 27.1.0.2     & 27.0.0.1 / --- \\
    sistemas  & Student S1        & 27.100.0.34  & 27.0.0.1 / 27.1.0.2 \\
    sistemas  & Student S2        & 27.100.0.35  & 27.0.0.1 / 27.1.0.2 \\
    civil     & www.civil.com     & 27.200.0.53  & 27.0.0.1 / 27.1.0.2 \\
    civil     & mail.civil.com    & 27.200.0.52  & 27.0.0.1 / 27.1.0.2 \\
    civil     & Student C1        & 27.200.0.55  & 27.0.0.1 / 27.1.0.2 \\
    civil     & Student C2        & 27.200.0.54  & 27.0.0.1 / 27.1.0.2 \\
    electrica & Student E1        & 27.150.0.101 & 27.150.0.1 / 27.1.0.2 \\
    electrica & Mail              & 27.150.0.102 & 27.150.0.1 / 27.1.0.2 \\
    electrica & Student E2        & 27.150.0.103 & 27.150.0.1 / 27.1.0.2 \\
    \bottomrule
  \end{tabularx}
  \caption{Direccionamiento de la red simulada.}
\end{table}

\fig{fig-pt-ipconfig.png}{Configuración estática de PC2: IP \texttt{27.100.0.34}, máscara \texttt{255.0.0.0}, gateway \texttt{27.0.0.1} y DNS \texttt{27.1.0.2}.}

\subsubsection{Servicios}

\paragraph{DNS.} En el servidor \texttt{27.1.0.2} se cargaron los registros de los tres
dominios (registros \texttt{A} para los servidores de correo y web, y \texttt{CNAME} para
los alias \texttt{www}, \texttt{pop3} y \texttt{smtp}).

\fig{fig-pt-dns.png}{Registros DNS cargados en el servidor 27.1.0.2.}
\fig{fig-pt-nslookup.png}{Resolución de \texttt{smtp.sistemas.com} consultando al DNS \texttt{27.1.0.2}.}

\paragraph{HTTP.} En los servidores web se activó el servicio y se personalizó la página
\texttt{index.html} de cada facultad.

\fig{fig-pt-http.png}{Acceso a la web por IP (\texttt{27.100.0.32}).}
\fig{fig-pt-http-url.png}{Acceso a la web por URL (\texttt{www.sistemas.com}), lo que prueba el DNS y el alias CNAME.}

\paragraph{Correo.} En cada servidor de correo se configuró el dominio y las cuentas de
usuario (una por equipo cliente).

\fig{fig-pt-email-client.png}{Configuración del cliente de correo en PC2 (POP3/SMTP de \texttt{sistemas.com}).}
\fig{fig-pt-email-server.png}{Cuentas de correo creadas en el servidor de sistemas.}
\fig{fig-pt-email-recibido.png}{Recepción de un correo entre clientes del mismo dominio.}
\fig{fig-pt-email-entre-dominios.png}{Recepción de un correo entre dominios: de \texttt{PC0@civil.com} a \texttt{PC2@sistemas.com}.}

\paragraph{FTP.} Se habilitó el servicio FTP en el servidor web de sistemas y se creó un
usuario con permisos de lectura y escritura.

\fig{fig-pt-ftp-server.png}{Servicio FTP y usuarios configurados en el servidor de sistemas.}
\fig{fig-pt-ftp.png}{Transferencia de un archivo por FTP desde el cliente (comando \texttt{put}).}
\fig{fig-pt-ftp-get.png}{Listado de archivos (\texttt{dir}) y descarga (\texttt{get}) por FTP; se aprecia \texttt{datos.txt}, subido previamente.}

\subsubsection{Pruebas y simulación}

\fig{fig-pt-ping.png}{Prueba de conectividad con el servidor DNS (\texttt{ping 27.1.0.2}).}
\fig{fig-pt-simulacion.png}{Inspección de un PDU HTTP en modo Simulation: capa de aplicación (HTTP), puertos de transporte (80 y 1061) y direccionamiento IP/MAC.}

\subsection{Parte B --- Red real}

\subsubsection{Wireshark}

\textbf{Wireshark} es un analizador de protocolos: captura las tramas que entran y
salen de la tarjeta de red y las decodifica capa por capa. A diferencia de Packet
Tracer, aquí se observa el tráfico \textbf{real}, no simulado. Las capturas se
hicieron sobre la interfaz \textbf{Wi-Fi}, no sobre las interfaces virtuales de
VirtualBox ni sobre el túnel de Tailscale, donde no circula el tráfico de la red.

Un punto clave de toda la actividad: la capa de aplicación solo se ve si el
protocolo viaja \textbf{sin cifrar}. Sobre \texttt{https://} el HTTP va dentro de
TLS y Wireshark solo muestra metadatos; por eso las pruebas se hacen contra sitios
\texttt{http://}, y ese contraste es en sí mismo uno de los resultados del
laboratorio.

\paragraph{Consulta web.} Se capturó mientras se abría el sitio del laboratorio de
computadores (\texttt{laboratorio.is.escuelaing.edu.co}). En la lista se distinguen
las \textbf{peticiones} (origen \texttt{192.168.1.11}, el equipo propio) de las
\textbf{respuestas} (origen \texttt{45.239.88.92}, el servidor), y se identifican los
protocolos de capa de aplicación activos: \textbf{HTTP} sobre el puerto
\textbf{80/TCP}, y \textbf{DNS} sobre \textbf{53/UDP} en la resolución previa del
nombre. En el detalle de la petición se leen los campos en texto plano
---\texttt{GET / HTTP/1.1}, \texttt{Host}, \texttt{User-Agent}, \texttt{Accept},
\texttt{Accept-Language} y la cabecera \texttt{Cookie}--- y en la respuesta,
\texttt{HTTP/1.1 200 OK} con \texttt{Content-Type: text/html}.

\fig{fig-wireshark-web.png}{Captura de una consulta web al sitio del laboratorio: la petición \texttt{GET / HTTP/1.1} (paquete 118) y su respuesta \texttt{200 OK} (paquete 123), con la capa de aplicación desplegada.}

Cuatro observaciones que salen de esta captura y alimentan el análisis:

\begin{itemize}
  \item El sitio es una \textbf{aplicación de una sola página}: después del documento
        inicial el navegador pide \texttt{assets/index-*.js} y \texttt{index-*.css}, y
        luego cada imagen por separado. Por eso una sola visita genera decenas de
        peticiones.
  \item Varias imágenes responden \texttt{404 Not Found}
        (\texttt{/images/labs/portada2.jpg}, por ejemplo). El servidor contesta
        correctamente el protocolo, pero el recurso no existe: son enlaces muertos que
        quedaron en el sitio.
  \item La cabecera \texttt{Cookie} viaja \textbf{en texto plano}: se ven
        \texttt{\_gcl\_au} y \texttt{\_ga}, cookies de seguimiento publicitario. Es la
        mejor demostración de por qué \texttt{HTTPS} es necesario: sobre HTTP,
        cualquier equipo del camino puede leer esa información.
  \item El destino \texttt{45.239.88.92} pertenece al bloque \texttt{45.239.88.0/22}
        de la Escuela que se determinó en la sección de WHOIS: el sitio del
        laboratorio está alojado en la red institucional.
\end{itemize}

\paragraph{Tráfico DHCP.} Se capturó el intercambio completo provocado por
\texttt{ipconfig /release} seguido de \texttt{ipconfig /renew}. Los mensajes se
identifican por su \emph{Transaction ID}: primero un \texttt{Release} (que libera la
concesión) y después la secuencia \textbf{DORA} completa ---\emph{Discover},
\emph{Offer}, \emph{Request} y \emph{Ack}--- que comparten el mismo identificador de
transacción, y por eso se agrupan como una sola conversación.

En la capa de transporte, DHCP usa siempre los mismos puertos: el cliente consulta
desde \textbf{68/UDP} al \textbf{67/UDP} del servidor. Y hay un detalle que confunde
a primera vista: en el \emph{Discover} el origen es \texttt{0.0.0.0} y el destino
\texttt{255.255.255.255} (broadcast), porque el equipo \textbf{todavía no tiene
dirección}: se identifica solo con su MAC y pregunta al broadcast si hay algún
servidor de direcciones disponible. Lo mismo ocurre con el \emph{Request}. Recién en
el \emph{Ack} el servidor confirma la dirección asignada.

En la capa de aplicación, el \emph{Discover} lleva el tipo de mensaje (\emph{Boot
Request}), la MAC del cliente y las opciones \textbf{53} (tipo de mensaje DHCP),
\textbf{61} (identificador del cliente), \textbf{12} (nombre del equipo),
\textbf{60} (clase de proveedor) y \textbf{55} (lista de parámetros solicitados).
Esa última es la que le dice al servidor \emph{qué} información extra quiere el
cliente además de la IP: router, servidores DNS, máscara y tiempo de concesión.

\fig{fig-wireshark-dhcp.png}{Intercambio DHCP completo: el \texttt{Release} y la secuencia \emph{Discover}, \emph{Offer}, \emph{Request} y \emph{Ack}, con los puertos \texttt{68} y \texttt{67} del \emph{Discover} a la vista.}

\paragraph{TELNET vs HTTP.} Para cerrar la actividad se repitió la consulta
\textbf{sin navegador}: se envió una petición \texttt{GET} escrita a mano contra el
puerto 80 del servidor, que es exactamente lo que el navegador hace por debajo. Se
utilizó un cliente TCP crudo equivalente a TELNET, porque el cliente TELNET de
Windows inyectaba bytes de negociación de opciones antes de la petición y el
servidor los interpretaba como una petición malformada, respondiendo
\texttt{400 Bad Request}. Esa conexión fallida igual sirvió: su ciclo de vida
completo (\texttt{SYN}, \texttt{SYN/ACK}, \texttt{ACK}, datos, \texttt{FIN}) quedó
capturado y permitió localizar el problema.

La petición enviada fue:

\begin{verbatim}
GET /~csantiago/RECO/index.html HTTP/1.0
Host: profesores.is.escuelaing.edu.co
\end{verbatim}

y el servidor respondió \texttt{200 OK}, con los headers
\texttt{Server: Apache/2.4.53 (Unix) PHP/8.1.4}, \texttt{Content-Type: text/html} y
\texttt{Content-Length: 242}, seguidos del \textbf{HTML en texto plano}. La ventana
\emph{Follow TCP Stream} de Wireshark permite ver la conversación completa, con la
petición del cliente y la respuesta del servidor diferenciadas por color.

La diferencia con el navegador es entonces evidente: el navegador envía la misma
petición y recibe la misma respuesta, pero además \textbf{interpreta el HTML y lo
renderiza}, y descarga los binarios como archivos. Con un cliente de texto, en
cambio, el contenido llega \textbf{crudo}: al pedir \texttt{prueba.pdf} o
\texttt{network.png}, lo que aparece en pantalla son los bytes del archivo (la firma
\texttt{\%PDF-1.5} o la de PNG) y basura ilegible, porque no hay nada que los
interprete. El aviso \emph{Not secure} que muestra el navegador confirma, además,
que el sitio viaja por HTTP sin cifrado.

\fig[0.85]{fig-wireshark-telnet-http.png}{Conversación completa capturada en Wireshark: la petición \texttt{GET} del cliente y la respuesta \texttt{200 OK} del servidor, ambas en texto plano.}
\fig[0.62]{fig-http-navegador.png}{La misma página cargada en el navegador: idéntico contenido, pero interpretado y renderizado. El aviso \emph{Not secure} confirma que la conexión es HTTP sin cifrar.}

\subsubsection{Consulta WHOIS y DNS (centralops.net)}

\textbf{Central Ops} es una utilidad web que reúne, en una sola consulta
(\emph{Domain Dossier}), tres registros distintos: el \textbf{WHOIS del dominio},
sus \textbf{registros DNS} y el \textbf{WHOIS de red} asociado a la IP. Conviene
distinguirlos desde el principio:

\begin{itemize}
  \item El \textbf{WHOIS del dominio} lo sirve el \emph{registro} del TLD
        (\texttt{whois.registry.co} para \texttt{.co}, \texttt{whois.internic.net}
        para \texttt{.com}) y responde quién registró el dominio, cuándo, con quién
        y hasta cuándo.
  \item El \textbf{WHOIS de red} lo sirve el \emph{registro regional de Internet}
        (RIR) dueño del bloque de direcciones (LACNIC, ARIN, RIPE) y responde qué
        rango IP está asignado, a qué organización y desde cuándo.
\end{itemize}

\paragraph{Limitación encontrada en los dominios \texttt{.co}.} La consulta del
WHOIS de dominio para \texttt{escuelaing.edu.co} y \texttt{jbb.gov.co}
\textbf{falla} en Central Ops con el error \texttt{CouldNotResolveServerName}. La
causa no es la red local: la herramienta tiene fijado el servidor
\texttt{whois.nic.co}, que \textbf{ya no existe} (su nombre no resuelve por DNS);
el registro colombiano migró el servicio a \texttt{whois.registry.co}, que sí
resuelve. El WHOIS de red y los registros DNS, en cambio, funcionan con
normalidad. Para completar los datos de dominio de esos dos casos se consultó el
registro del propio registry. Este hallazgo se documenta a propósito: demuestra
que la utilidad web es una comodidad de presentación y que \textbf{la fuente de
verdad es siempre el registro}.

\paragraph{Dominio \texttt{escuelaing.edu.co}.}

\begin{table}[H]
  \centering
  \begin{tabularx}{\linewidth}{@{}lX@{}}
    \toprule
    \textbf{Consulta} & \textbf{Resultado} \\
    \midrule
    Servidores de dominio (NS) & \textbf{2}: \texttt{ns1.escuelaing.edu.co} y \texttt{ns2.escuelaing.edu.co} \\
    Asignado & 2 de junio de 1998 (hace 28 años) \\
    Registrado con & \texttt{www.registrocolombia.co} (Registro de Colombia) \\
    ID de la entidad registradora & \texttt{111111} \\
    Última actualización & 19 de julio de 2025 \\
    Vigencia del registro & hasta el 31 de diciembre de 2028 \\
    Rango IP y autoridad & \texttt{45.239.88.0/22} --- \textbf{LACNIC} \\
    Empresa asignada & Escuela Colombiana de Ingeniería (\texttt{CO-ECIN2-LACNIC}) \\
    \bottomrule
  \end{tabularx}
  \caption{Resultados de la consulta WHOIS/DNS de \texttt{escuelaing.edu.co}.}
\end{table}

El dato más revelador está en el registro \texttt{SOA}: su servidor primario es
\texttt{ns1.escuelaing.edu.co} y su contacto \texttt{admin@escuelaing.edu.co}, lo
que significa que \textbf{la Escuela hospeda su propio DNS}. Su bloque
\texttt{45.239.88.0/22} está asignado como \emph{miembro directo} de LACNIC
(identificador \texttt{CO-ECIN2-LACNIC}), no por intermedio de un proveedor. El
correo, en cambio, \emph{sí} está en la nube: el registro \texttt{MX} apunta a
\texttt{escuelaing-edu-co.mail.protection.outlook.com} (Microsoft 365) y la
política \texttt{SPF} autoriza, además de Outlook, a Brevo
(\texttt{\_spf.embluemail.com}) y HubSpot para el envío de campañas. Los registros
\texttt{TXT} de verificación dejan ver los servicios externos integrados (Google,
Facebook, GlobalSign y Sophos).

\fig[0.72]{fig-centralops-escuelaing-whois.png}{Consulta de \texttt{escuelaing.edu.co}: el WHOIS de dominio falla en Central Ops mientras el WHOIS de red (LACNIC) responde con el bloque \texttt{45.239.88.0/22} de la Escuela.}
\fig[0.95]{fig-centralops-escuelaing-dns.png}{Registros DNS de \texttt{escuelaing.edu.co}: \texttt{A}, \texttt{MX} (Microsoft 365), \texttt{TXT} (SPF y verificaciones), \texttt{SOA} y los dos \texttt{NS}.}

\paragraph{Dominio \texttt{jbb.gov.co}.}

\begin{table}[H]
  \centering
  \begin{tabularx}{\linewidth}{@{}lX@{}}
    \toprule
    \textbf{Consulta} & \textbf{Resultado} \\
    \midrule
    Servidores de dominio (NS) & \textbf{2}: \texttt{rick.ns.cloudflare.com} y \texttt{tori.ns.cloudflare.com} \\
    Asignado & 20 de enero de 2000 (hace 26 años) \\
    Registrado con & \texttt{www.registrocolombia.co} (Registro de Colombia) \\
    ID de la entidad registradora & \texttt{111111} \\
    Última actualización & 14 de mayo de 2021 \\
    Vigencia del registro & hasta el 20 de enero de 2030 \\
    Rango IP y autoridad & \texttt{20.64.0.0/10} --- \textbf{ARIN} \\
    Empresa asignada & Microsoft Corporation (\texttt{MSFT}) \\
    \bottomrule
  \end{tabularx}
  \caption{Resultados de la consulta WHOIS/DNS de \texttt{jbb.gov.co}.}
\end{table}

Este caso es el \textbf{opuesto} al anterior en cuanto a infraestructura. Sus
servidores de nombre están en \textbf{Cloudflare} y las direcciones a las que
resuelve (\texttt{20.94.123.146} y \texttt{20.119.228.39}) pertenecen a
\textbf{Microsoft Corporation}: el bloque \texttt{20.64.0.0/10} que las contiene
está registrado en ARIN como \texttt{MSFT}. Es decir, el organismo \textbf{delegó
tanto el DNS como el alojamiento a terceros en la nube}, aunque el dominio siga
registrado en Colombia. Los registros inversos lo confirman al mostrar los
servidores de \textbf{Azure DNS} (\texttt{ns1-07.azure-dns.com},
\texttt{ns2-07.azure-dns.net}). La lección de fondo: \emph{el dueño del dominio,
quien sirve el DNS y quien aloja el servicio pueden ser tres organizaciones
distintas}.

\fig[0.65]{fig-centralops-jbb-whois.png}{Consulta de \texttt{jbb.gov.co}: el WHOIS de dominio vuelve a fallar por la limitación de \texttt{.co}, pero el WHOIS de red (ARIN) identifica el bloque como propiedad de Microsoft Corporation.}
\fig[0.80]{fig-centralops-jbb-dns.png}{Registros DNS de \texttt{jbb.gov.co}: los dos \texttt{NS} están en Cloudflare y las consultas inversas revelan Azure DNS.}

\paragraph{Dominio \texttt{google.com}.}

\begin{table}[H]
  \centering
  \begin{tabularx}{\linewidth}{@{}lX@{}}
    \toprule
    \textbf{Consulta} & \textbf{Resultado} \\
    \midrule
    Servidores de dominio (NS) & \textbf{4}: \texttt{ns1}, \texttt{ns2}, \texttt{ns3} y \texttt{ns4.google.com} \\
    Asignado & 15 de septiembre de 1997 (hace 29 años) \\
    Registrado con & MarkMonitor Inc. (registrador corporativo privado) \\
    ID de la entidad registradora & \texttt{292} \\
    Última actualización & 9 de septiembre de 2019 \\
    Vigencia del registro & hasta el 14 de septiembre de 2028 \\
    Rango IP y autoridad & \texttt{172.253.0.0/16} --- \textbf{ARIN} \\
    Empresa asignada & Google LLC (\texttt{GOGL}) \\
    \bottomrule
  \end{tabularx}
  \caption{Resultados de la consulta WHOIS/DNS de \texttt{google.com}.}
\end{table}

En \texttt{google.com} todo es corporativo y propio: \textbf{cuatro} servidores de
nombre bajo su propio dominio, un bloque \texttt{172.253.0.0/16} asignado
\emph{directamente} a Google LLC por ARIN, y el registro del dominio en manos de
\textbf{MarkMonitor}, un registrador privado que gestiona carteras de marcas. Los
seis estados del dominio (\texttt{client} y \texttt{server} \texttt{Delete},
\texttt{Transfer} y \texttt{Update} \emph{Prohibited}) son un mecanismo
anti-secuestro: impiden que el dominio se transfiera o se borre, incluso por error
del propio dueño. La política \texttt{SPF} también es restrictiva, coherente con
una empresa que solo envía correo desde su propia infraestructura.

\fig[0.66]{fig-centralops-google-whois.png}{Consulta de \texttt{google.com}: el WHOIS de dominio sí funciona en Central Ops (el TLD \texttt{.com} usa Verisign) y muestra el registrador MarkMonitor y los cuatro servidores de nombre.}
\fig[0.72]{fig-centralops-google-neto.png}{WHOIS de red de \texttt{google.com} en ARIN: el bloque \texttt{172.253.0.0/16} está asignado directamente a Google LLC.}
\fig[0.62]{fig-centralops-google-dns.png}{Registros DNS de \texttt{google.com}: seis direcciones \texttt{A}, cuatro \texttt{AAAA}, los cuatro \texttt{NS}, el \texttt{MX} y un \texttt{SOA} con valores de refresco mucho más agresivos (900~s) que los de la Escuela (86\,400~s).}

\paragraph{Dominio \texttt{cern.ch} (organización no americana).}

\begin{table}[H]
  \centering
  \begin{tabularx}{\linewidth}{@{}lX@{}}
    \toprule
    \textbf{Consulta} & \textbf{Resultado} \\
    \midrule
    Servidores de dominio (NS) & \textbf{3}: \texttt{ext-dns-1.cern.ch}, \texttt{ext-dns-2.cern.ch} y \texttt{scsnms.switch.ch} \\
    Asignado & \emph{no publicado} por el registro \texttt{.ch} (SWITCH) \\
    Registrado con & CSC Corporate Domains, Inc. (Wilmington, EE. UU.) \\
    ID de la entidad registradora & \emph{no publicado} por el registro \texttt{.ch} \\
    Última actualización & \emph{no publicado} por el registro \texttt{.ch} \\
    Vigencia del registro & \emph{no publicado} por el registro \texttt{.ch} \\
    Rango IP y autoridad & \texttt{188.184.0.0/16} --- \textbf{RIPE} \\
    Empresa asignada & CERN --- European Organization for Nuclear Research (\texttt{ORG-CE0-RIPE}, AS513) \\
    \bottomrule
  \end{tabularx}
  \caption{Resultados de la consulta WHOIS/DNS de \texttt{cern.ch}.}
\end{table}

\paragraph{Segunda limitación: el registro \texttt{.ch} no publica fechas.} A
diferencia del registro colombiano, SWITCH \textbf{no expone} la fecha de creación,
la de expiración ni el identificador del registrador: su respuesta RDAP trae el
campo de eventos \textbf{vacío} y su servidor WHOIS rechaza las consultas
automatizadas. Por eso cuatro de las ocho preguntas del laboratorio \textbf{no
tienen respuesta pública} para este dominio. Se documenta como hallazgo:
\emph{los registros publican distinta cantidad de información}, y saber cuál
publica qué es parte del trabajo.

Del bloque de red, en cambio, la información está completa: el rango
\texttt{188.184.0.0/16} está asignado al CERN por \textbf{RIPE NCC} como miembro
directo (\texttt{ORG-CE0-RIPE}), con el sistema autónomo \texttt{AS513} y contacto
\texttt{ext@p.cern.ch}. Es el tercer contraste del laboratorio: un organismo de
investigación europeo, con su dominio registrado en manos de una empresa
estadounidense (CSC) y su bloque IP administrado por el RIR europeo.

\fig[0.60]{fig-centralops-cern-whois.png}{Consulta de \texttt{cern.ch}: el WHOIS de dominio se resuelve por RDAP y muestra el registrador CSC, pero \textbf{sin fechas}; el WHOIS de red (RIPE) identifica el bloque \texttt{188.184.0.0/16} como del CERN.}
\fig[0.85]{fig-centralops-cern-dns.png}{Registros DNS de \texttt{cern.ch}: el \texttt{SOA} y los tres \texttt{NS} (dos propios del CERN y uno de SWITCH), con el \texttt{AAAA} del sitio y la resolución inversa.}

\paragraph{Comparación de los cuatro dominios.}

\begin{table}[H]
  \centering
  \footnotesize
  \begin{tabularx}{\linewidth}{@{}lXXXX@{}}
    \toprule
    & \texttt{escuelaing.edu.co} & \texttt{jbb.gov.co} & \texttt{google.com} & \texttt{cern.ch} \\
    \midrule
    Servidores DNS & 2 (propios) & 2 (Cloudflare) & 4 (propios) & 3 (CERN + SWITCH) \\
    Asignado & 02/06/1998 & 20/01/2000 & 15/09/1997 & no publicado \\
    Registrador & registrocolombia.co & registrocolombia.co & MarkMonitor Inc. & CSC Corporate Domains \\
    IANA ID & 111111 & 111111 & 292 & no publicado \\
    Vigente hasta & 31/12/2028 & 20/01/2030 & 14/09/2028 & no publicado \\
    Rango IP & 45.239.88.0/22 & 20.64.0.0/10 & 172.253.0.0/16 & 188.184.0.0/16 \\
    Autoridad & LACNIC & ARIN & ARIN & RIPE \\
    Empresa & Escuela Colombiana de Ingeniería & Microsoft Corporation & Google LLC & CERN \\
    \bottomrule
  \end{tabularx}
  \caption{Comparación de los cuatro dominios consultados.}
\end{table}

Los dos dominios colombianos comparten registrador (el Registro de Colombia) y su
identificador \texttt{111111} es un valor genérico del registry, no el ID real del
registrador. Sin embargo, su infraestructura es opuesta: la Escuela mantiene el DNS
\emph{on-premise} y su propio bloque IP ante LACNIC, mientras que JBB delegó todo a
Cloudflare y Microsoft. El dominio corporativo, por su parte, muestra el patrón de
una multinacional: registrador privado, DNS y red propios, y bloqueos anti-secuestro
activos. El dominio europeo de investigación cierra el cuadro: su red la administra
RIPE y su registro es estadounidense, recordatorio de que la sede del registrador no
tiene por qué coincidir con la del organismo.

\subsubsection{Servidor NTP}

\paragraph{Por qué importa que todos los equipos tengan la misma hora.} En una
infraestructura, la hora no es un detalle cosmético: es una \textbf{referencia
común} sin la cual muchos servicios dejan de funcionar o pierden sentido. Los
casos más claros son:

\begin{itemize}
  \item \textbf{Correlación de eventos y registros.} Para reconstruir qué pasó
        primero cuando varias máquinas registran un mismo incidente, los
        \emph{logs} tienen que estar en la misma escala de tiempo. Sin hora común,
        el orden de los hechos es indemostrable.
  \item \textbf{Autenticación centralizada.} Los protocolos como Kerberos
        (base de Active Directory) rechazan las peticiones cuyo desfase supera
        unos pocos minutos.
  \item \textbf{Certificados y seguridad.} La validez de un certificado se evalúa
        por fechas: un reloj adelantado o atrasado puede hacer que un certificado
        válido se vea vencido.
  \item \textbf{Servicios de red.} Correo, bases de datos replicadas y
        transacciones distribuidas dependen de marcas de tiempo consistentes.
  \item \textbf{Auditoría y forense.} Toda investigación posterior se apoya en la
        cronología de los registros.
\end{itemize}

\paragraph{Diseño del laboratorio.} Se siguió el reparto sugerido para grupos de
dos estudiantes: \textbf{Solaris actúa como servidor NTP} y las demás máquinas
(\textbf{Slackware}, \textbf{Windows Server con GUI} y \textbf{Windows Server
Core}) son clientes. Todas se comunican por la red interna virtual \texttt{lab-uni}
(\texttt{192.168.82.0/24}), que es la misma que se usó en el laboratorio de DNS.

Como esa red es \textbf{aislada y sin salida a Internet}, el servidor no puede
sincronizarse con una fuente externa. En ese escenario se usa el \textbf{reloj
local} de la propia máquina como referencia (\emph{driver} 1 de NTP), con un
estrato alto para dejar claro que no es una fuente autorizada:

\begin{verbatim}
server 127.127.1.0
fudge  127.127.1.0 stratum 8

restrict default nomodify notrap nopeer
restrict 127.0.0.1
restrict 192.168.82.0 mask 255.255.255.0 nomodify notrap
\end{verbatim}

Las líneas \texttt{restrict} dejan que los clientes \textbf{consulten} la hora,
pero impiden que la modifiquen o que se establezcan como pares del servidor. En
Solaris el servicio se administra por SMF y se habilita con
\texttt{svcadm enable network/ntp}.

En los clientes, la configuración apunta siempre al servidor:

\begin{itemize}
  \item \textbf{Slackware:} se reemplazó \texttt{/etc/ntp.conf} por una única línea
        \texttt{server 192.168.82.11 iburst} y se arrancó el demonio con
        \texttt{/etc/rc.d/rc.ntpd start}.
  \item \textbf{Windows (GUI y Core):} se configuró el servicio de hora con
        \texttt{w32tm /config /manualpeerlist:"192.168.82.11,0x8"
        /syncfromflags:manual /update}, se reinició el servicio y se forzó la
        sincronización con \texttt{w32tm /resync}.
\end{itemize}

\paragraph{Resultados.}

\begin{table}[H]
  \centering
  \begin{tabularx}{\linewidth}{@{}llX@{}}
    \toprule
    \textbf{Máquina} & \textbf{Rol} & \textbf{Estado verificado} \\
    \midrule
    Solaris \texttt{192.168.82.11} & Servidor & Servicio SMF \texttt{online}; \texttt{ntpq -p} muestra \texttt{LOCAL(0)} como referencia (estrato 8) y el servidor responde en el estrato 9 \\
    Slackware \texttt{192.168.82.10} & Cliente & \texttt{ntpd} en ejecución apuntando a \texttt{192.168.82.11} \\
    Windows GUI \texttt{192.168.82.13} & Cliente & \texttt{Stratum 10}, \texttt{ReferenceId 0xC0A8520B} (= 192.168.82.11), \texttt{Leap Indicator: 0} \\
    Windows Core \texttt{192.168.82.12} & Cliente & \texttt{Stratum 10}, \texttt{ReferenceId 0xC0A8520B} (= 192.168.82.11) \\
    \bottomrule
  \end{tabularx}
  \caption{Estado del servidor y de los clientes NTP al cierre de la actividad.}
\end{table}

\fig[0.85]{fig-ntp-server.png}{Servidor NTP en Solaris: el servicio \texttt{svc:/network/ntp} está \texttt{online} y \texttt{ntpq -p} muestra el reloj local como referencia en el estrato 8, con el registro de alcance completo (\texttt{reach 377}).}
\fig[0.90]{fig-ntp-client.png}{Cliente Windows sincronizado con el servidor: \texttt{w32tm} informa estrato 10 y el \texttt{ReferenceId 0xC0A8520B}, que corresponde a la dirección del servidor (\texttt{192.168.82.11}).}

\paragraph{Observaciones de la práctica.} Más allá del resultado, aparecieron
varios comportamientos que vale la pena registrar porque son los que ocurren en un
entorno virtualizado:

\begin{itemize}
  \item \textbf{La sincronización no es inmediata.} El servidor tardó varios minutos
        en seleccionar el reloj local como referencia: al principio informaba
        \texttt{stratum=16} y \texttt{sync\_unspec} (sin sincronizar) y recién
        después de unos ocho minutos pasó a \texttt{leap\_none} y
        \texttt{clock\_sync}.
  \item \textbf{El cliente no siempre ajusta de golpe.} NTP corrige con
        \emph{slew} (ajuste gradual acotado), no con saltos instantáneos. Por eso
        un desfase de varios segundos puede tardar en desaparecer; el ajuste
        puntual se logra con una consulta única de \texttt{ntpdate}.
  \item \textbf{El script de arranque de Slackware no era ejecutable.} El archivo
        \texttt{/etc/rc.d/rc.ntpd} existía pero sin el bit de ejecución
        (\texttt{-rw-r--r--}), por lo que el demonio no arrancaba en el inicio. Se
        corrigió con \texttt{chmod 755}.
  \item \textbf{Windows puede quedar en un estado que no acepta los datos.} El
        servicio \texttt{w32time} de la máquina con GUI respondía a las consultas
        NTP del servidor pero no las aplicaba; fue necesario
        \texttt{w32tm /unregister} seguido de \texttt{w32tm /register} para
        restablecerlo.
  \item \textbf{El hipervisor compite con NTP.} VirtualBox sincroniza el reloj de
        las máquinas virtuales con el del equipo anfitrión. En la VM de Slackware
        ese mecanismo y \texttt{ntpd} tiraban en direcciones distintas, y se
        observó un \emph{jitter} muy alto (varios segundos) que impedía el ajuste
        fino. Es un efecto propio de la virtualización, no de la configuración
        NTP, y en máquinas físicas no aparece.
\end{itemize}

\subsubsection{Cableado estructurado}

\paragraph{Construcción de patch cords.} Se armaron dos cables RJ-45 siguiendo el
estándar de asignación de colores del par trenzado. Un cable \textbf{directo} usa
la misma norma en los dos extremos (\textbf{T568B}--T568B) y sirve para unir
dispositivos de \emph{distinto} tipo, como un equipo y un switch. Un cable
\textbf{cruzado} usa \textbf{T568A} en un extremo y \textbf{T568B} en el otro, y
se usa entre dispositivos del \emph{mismo} tipo (equipo--equipo o
switch--switch). El procedimiento en ambos casos es el mismo: pelar el recubrimiento
exterior, destrenzar los pares, ordenarlos según la norma, cortarlos parejos,
insertarlos en el conector RJ-45 y fijarlos con la ponchadora. La comprobación se
hace con el \textbf{probador de cables}.

\begin{verbatim}
T568B: 1 blanco-naranja  2 naranja   3 blanco-verde  4 azul
       5 blanco-azul     6 verde     7 blanco-marron 8 marron

T568A: 1 blanco-verde    2 verde     3 blanco-naranja 4 azul
       5 blanco-azul     6 naranja   7 blanco-marron  8 marron
\end{verbatim}

\fig[0.80]{fig-cable-referencia.jpg}{Referencia del orden de colores del cable RJ-45 tipo A y B, consultada durante el armado de los patch cords.}

\paragraph{Crimpado en patch panel y faceplates.} En el grupo de trabajo se realizó
la prueba de crimpado del \emph{cableado horizontal}: se conectaron dos equipos a
través de un \textbf{patch panel} y dos \textbf{faceplates}, cada uno con su toma de
información. La continuidad se verificó de las dos formas posibles: con el
probador de cables y con un \texttt{ping} entre los equipos conectados.

\fig[0.62]{fig-patch-panel.jpg}{Patch panel con los latiguillos ya crimpados y el faceplate con su toma de información.}
\fig[0.95]{fig-patch-panel-panorama.jpg}{Puesto de trabajo durante la prueba de crimpado: patch panel, faceplates, ponchadora y cable UTP.}

Como en ese momento los equipos estaban conectados directamente y sin un servidor
DHCP que repartiera direcciones, se autoasignaron direcciones \emph{link-local}
(\texttt{169.254.x.x}). Que el \texttt{ping} entre ellos respondiera demuestra que
la continuidad eléctrica del cableado era correcta.

\fig[0.95]{fig-cable-ping.jpg}{Prueba de continuidad: \texttt{ping} exitoso entre los dos equipos unidos por el patch panel (direcciones link-local \texttt{169.254.x.x}, asignadas automáticamente al no haber DHCP).}

\paragraph{Patch cords directo y cruzado.} Los dos latiguillos se armaron sobre la
mesa de trabajo con cable UTP y conectores RJ-45: uno \textbf{directo}, con la norma
T568B en los dos extremos, y uno \textbf{cruzado}, con T568A en un extremo y T568B en
el otro. En ambos casos el procedimiento fue el mismo: pelar el recubrimiento
exterior, destrenzar los pares, ordenarlos según la norma, cortarlos parejos,
insertarlos en el conector y fijarlos con la ponchadora. La bolsa de conectores, con
el diagrama de colores impreso, quedó a la vista sobre la mesa durante el armado.

\fig[0.92]{fig-cable-cords.jpg}{Latiguillos ya crimpados y conectados al patch panel durante el armado de los patch cords. Se distinguen los conectores RJ-45 ponchados y el cable UTP tendido sobre la mesa de trabajo.}

La continuidad de cada latiguillo se comprobó con el probador de cables y, además,
quedó respaldada por la prueba de conectividad entre los dos equipos descrita más
arriba: ese \texttt{ping} sin pérdidas recorre justamente la cadena
latiguillo--patch panel--cable horizontal--faceplate--latiguillo, de modo que
también sirve como comprobación de los dos cables armados.

\paragraph{Cableado estructurado del edificio.} El cableado estructurado organiza la
instalación en subsistemas, y el objetivo del punto es reconocerlos sobre la
instalación real. El \textbf{cableado horizontal} es el tramo fijo que va de cada toma
de información al cuarto de comunicaciones; el \textbf{patch panel} es el punto donde
ese tendido queda terminado, y cumple del lado del cuarto la misma función que el
\textbf{faceplate} cumple del lado del puesto: fijar la terminación para que el cable
fijo no se mueva cada vez que se cambia una conexión. Los \textbf{latiguillos} (patch
cords) son los tramos móviles del esquema: uno une el equipo con la toma del faceplate
y otro une el puerto del patch panel con el switch. El tendido se conduce por
\textbf{canaletas}, que lo protegen mecánicamente y mantienen el orden de la
instalación.

La ventaja del diseño es que separa lo fijo de lo móvil: para reorganizar una conexión
se cambia un latiguillo, no se recablea el tendido del edificio. En la práctica se
trabajó con los elementos de terminación de ese esquema --patch panel, faceplates y
latiguillos--, que son los mismos que aparecen concentrados en el rack del cuarto de
comunicaciones, como muestran las figuras anteriores. El registro fotográfico de esta
sección corresponde a los elementos de terminación con los que se practicó; no se
incluye registro del cuarto de comunicaciones del edificio.

\section{Conclusiones}

El laboratorio permitió ver, sobre tráfico y equipos reales, los protocolos que en
los laboratorios anteriores se habían estudiado en simulación. Las conclusiones
principales son:

\begin{itemize}
  \item \textbf{La capa de aplicación solo es visible cuando no hay cifrado.} Con
        Wireshark se leyeron en texto plano la petición \texttt{GET} y los headers
        de la respuesta, y se distinguieron los puertos de cada protocolo (80/TCP
        para HTTP, 53/UDP para DNS, 67 y 68/UDP para DHCP). El contraste fue
        directo: en el mismo ejercicio se comprobó que sobre \texttt{https://} el
        contenido viaja cifrado y solo se ven metadatos. La demostración más
        contundente fue ver la cabecera \texttt{Cookie} de seguimiento publicitario
        viajando sin cifrar.

  \item \textbf{El navegador es una capa de interpretación, no de transporte.} Al
        pedir la misma página con un cliente de texto se recibió exactamente el
        mismo contenido, pero \textbf{crudo}: el HTML sin renderizar y, al pedir un
        binario, la secuencia de bytes del archivo. El navegador hace por encima lo
        que el cliente de texto no hace: interpretar, dibujar y guardar.

  \item \textbf{El formato de las tramas no depende del contenido.} El ciclo de vida
        de una conexión TCP ---saludo de tres vías, transferencia y cierre--- y los
        mensajes del intercambio DHCP (los cuatro del proceso DORA) se observaron
        tal como los describe la teoría.

  \item \textbf{Los registros publican distinta cantidad de información.} La
        consulta WHOIS dejó tres lecciones concretas: la herramienta web puede
        quedar obsoleta (Central Ops apuntaba a un servidor WHOIS que ya no
        existe), cada registro decide qué publica (el registro \texttt{.ch} no
        expone fechas de creación ni de expiración) y \textbf{la fuente de verdad
        es siempre el registro}, no la utilidad que lo muestra. También se
        distinguió con claridad el WHOIS del dominio (lo sirve el registro del TLD)
        del WHOIS de red (lo sirve el RIR dueño del bloque IP: LACNIC, ARIN o
        RIPE), y se comprobó que el dueño del dominio, quien sirve el DNS y quien
        aloja el servicio pueden ser tres organizaciones distintas.

  \item \textbf{La hora es infraestructura.} Montar un servidor NTP y sincronizar
        los clientes mostró que la hora no es un dato accesorio: es la referencia
        común que hace posibles los registros correlacionables, la autenticación,
        la validez de los certificados y las transacciones. La jerarquía de
        estratos y la diferencia entre un servidor con referencia externa y uno
        que solo sirve su reloj local quedaron claras.

  \item \textbf{La virtualización tiene efectos propios.} Sincronizar la hora
        dentro de máquinas virtuales resultó menos limpio que en equipos físicos:
        el hipervisor también ajusta el reloj de la máquina virtual y compite con
        NTP, lo que se manifestó como un \emph{jitter} alto que impide el ajuste
        fino. No es un error de configuración, sino una característica del entorno,
        y conviene tenerla presente al interpretar resultados en laboratorio.

  \item \textbf{Las herramientas ayudan, pero no reemplazan el fundamento.} Tanto
        en WHOIS como en Wireshark, los obstáculos ---una URL con la tilde
        faltante, un cliente TELNET que inyecta caracteres de control, un registro
        que no publica sus datos, un servicio de hora que no arranca por un permiso
        de ejecución--- se resolvieron entendiendo \emph{qué} hace cada protocolo
        por debajo, no memorizando comandos. Ese es el aprendizaje que queda: la
        herramienta es una comodidad, y saber bajar un nivel es lo que permite
        diagnosticar.
\end{itemize}

\section{Bibliografía}

\begin{itemize}
  \item Cisco Networking Academy. \emph{Introduction to Networks}. Cisco Press.
  \item Documentación de Cisco Packet Tracer.
  \item Material de clase de Arquitectura y Servicios de Red (AYSR).
\end{itemize}

\end{document}
